\section{初版方案的独立测试}\label{sec:test}

\subsection{测试设置}

测试使用固定划分中的 32 张图，每张 $512\times512$，共 8,388,608 个有效像素。检测只读原图、训练统计和固定配置；评价程序再读取保存的预测，与独立生成的真值对齐。测试集没有参与选参，也没有在测试后修改配置。当前测试图中没有眼镜像素，忽略逻辑保留但未被触发。

评价只做像素计数，不调用检测函数内部的状态；核心循环见代码~\textbf{4-1}。

\noindent\textbf{代码 4-1}\quad 有效区域内的混淆计数与指标（evaluation.py）
\begin{lstlisting}
valid = (np.ones(gt.shape, dtype=bool) if valid_mask is None
         else valid_mask.astype(bool))
not_pred, not_gt = np.logical_not(pred), np.logical_not(gt)
tp = np.count_nonzero(pred & gt & valid)
fp = np.count_nonzero(pred & not_gt & valid)
tn = np.count_nonzero(not_pred & not_gt & valid)
fn = np.count_nonzero(not_pred & gt & valid)

precision = tp / (tp + fp) if tp + fp else None
recall    = tp / (tp + fn) if tp + fn else None
f1        = 2 * tp / (2 * tp + fp + fn) if (2 * tp + fp + fn) else None
iou       = tp / (tp + fp + fn) if (tp + fp + fn) else None
\end{lstlisting}

这段代码用逻辑与同时完成“预测为正”和“真值为正”的筛选：\texttt{pred} 是预测掩膜，\texttt{gt} 是脸部真值，二者都是 0/1 数组；\texttt{valid} 把眼镜像素排除在四个计数之外。分母为零时返回 \texttt{None}，不做静默填零。汇总时先把所有图像的 TP/FP/TN/FN 相加再计算 global，逐图平均则单独对每张图的指标做等权平均；两种口径在报告中始终分开标注。

\subsection{总体结果}

总体指标与混淆计数分别见表~\ref{tab:test} 和表~\ref{tab:testcounts}。

\begin{table}[H]
\centering\small
\caption{测试集总体指标，单位为百分比}\label{tab:test}
\begin{tabular}{lrrrr}
\toprule
方案 & Precision & Recall & F1 & IoU \\
\midrule
纯阈值基线 & 52.91 & 76.82 & 62.66 & 45.63 \\
初版方案 & 66.43 & 75.55 & 70.69 & 54.67 \\
变化，百分点 & +13.52 & $-1.27$ & +8.03 & +9.04 \\
\bottomrule
\end{tabular}
\end{table}

\begin{table}[H]
\centering\small
\caption{测试集混淆计数}\label{tab:testcounts}
\begin{tabular}{lrrrr}
\toprule
方案 & TP & FP & TN & FN \\
\midrule
纯阈值基线 & 1,652,227 & 1,470,496 & 4,767,390 & 498,495 \\
初版方案 & 1,624,818 & 821,271 & 5,416,615 & 525,904 \\
\bottomrule
\end{tabular}
\end{table}

初版相对基线：F1 提高 8.03 个百分点，IoU 提高 9.04 个百分点；Precision 提高 13.52，Recall 下降 1.27。计数上，FP 减少 649,225，FN 增加 27,409。收益主要来自减少误检，而不是找回漏检。

初版相对基线的变化主要是“删掉不该选的”，而不是“把漏掉的找回来”。召回基本不动，精确率大幅上升；这也解释了为什么初版的 F1 停在 70\% 附近——它把颜色候选修剪得更干净，但没有增加新的判别信息。

\subsection{提升来自哪一步}

把流程拆成四个阶段，可以看到每一步拿什么换什么，结果如表~\ref{tab:stage}。

\begin{table}[H]
\centering\small
\setlength{\tabcolsep}{4pt}
\caption{固定流程各阶段的测试结果，指标单位为百分比}\label{tab:stage}
\begin{tabular}{lrrrrr}
\toprule
阶段 & Precision & Recall & F1 & FP & FN \\
\midrule
基线阈值，$k=1.50$ & 52.91 & 76.82 & 62.66 & 1,470,496 & 498,495 \\
颜色阈值，$k=1.25$ & 56.59 & 65.24 & 60.61 & 1,076,315 & 747,523 \\
保留最大连通域 & 64.73 & 63.36 & 64.04 & 742,515 & 787,938 \\
填孔后的结果 & 66.43 & 75.55 & 70.69 & 821,271 & 525,904 \\
\bottomrule
\end{tabular}
\end{table}

三步的交换率不同。$k$ 从 1.50 收到 1.25，去掉 394,181 个误检像素，也丢掉 249,028 个正确像素，约合 1 个 TP 换 1.6 个 FP；最大连通域去掉 333,800 个误检，只损失 40,415 个正确像素，约合 1 个 TP 换 8.3 个 FP；填孔补回 262,034 个正确像素，同时增加 78,756 个误检，约合 1 个 FP 换 3.3 个 TP。

单看 $k$ 收紧这一步，F1 从 62.66 降到 60.61，是亏的；它是在给后面的最大连通域让路。最大连通域是最划算的一步；填孔再把内部补回来。三步不能分开评价：阈值、区域和填孔之间存在配合关系，单独调其中一步会得到不同的“最优”。

\subsection{误检从哪里来}

表~\ref{tab:errors} 把最终误检按原始类别归类，并给出填孔前后两组正类的组内召回。

\begin{table}[H]
\centering\small
\caption{最终误检的语义来源与五官召回}\label{tab:errors}
\begin{tabular}{lrrr}
\toprule
负类 & 基线 FP & 最终 FP & 占最终 FP \\
\midrule
头发 & 772,905 & 497,516 & 60.58\% \\
脖颈 & 248,769 & 192,139 & 23.40\% \\
背景 & 369,581 & 93,639 & 11.40\% \\
衣服 & 58,056 & 20,420 & 2.49\% \\
双耳 & 17,665 & 14,184 & 1.73\% \\
其他 & 3,520 & 3,373 & 0.41\% \\
\bottomrule
\end{tabular}

\vspace{6pt}
\begin{tabular}{lrr}
\toprule
正类分组 & 最大区域后 & 填孔后 \\
\midrule
皮肤与鼻子 & 66.43 & 76.27 \\
眼、眉、口腔与嘴唇 & 32.61 & 68.32 \\
\bottomrule
\end{tabular}
\end{table}

头发和脖颈合计占最终误检的 83.98\%。头发容易形成类肤色候选；脖颈则与脸部颜色相近，又常常连在同一前景块上，最大连通域无法把它切开。所以剩余问题不只来自分离的背景，也来自与脸部同色或相连的区域。

其中脖子和耳朵本身是皮肤，这是方法的语义上限。测试集中脖子与双耳合计约 37.4 万像素，而脸部正类约 215 万像素。即使颜色判决完全正确、不再产生任何头发或背景误检，只要结构仍是“最大连通域加填孔”，这部分像素仍会被计入前景，精确率的上限约为 85\%。当前精确率 66.43\%，离上限还有约 19 个百分点；即使把头发和背景误检全部消除，纯颜色模型也无法达到更高的精确率。这个计算假设所有肤色像素都被判为正类，实际方法可能因连通性排除部分耳朵，因此 85\% 是量级上的上限，不是严格界。

填孔确实恢复了五官：眼、眉、口腔与嘴唇的组内召回从 32.61\% 提高到 68.32\%，补回 69,662 个正确像素；皮肤与鼻子组从 66.43\% 提高到 76.27\%。方法并不知道那些像素是眼睛或嘴，它只是利用了它们被脸部轮廓包围这一拓扑事实。

\subsection{逐图稳定性与亮度条件}

逐图指标与亮度分组见表~\ref{tab:perimage}。

\begin{table}[H]
\centering\small
\caption{逐图指标、边界质量与亮度分组}\label{tab:perimage}
\begin{tabular}{lrr}
\toprule
统计项目 & 基线 & 初版 \\
\midrule
逐图 F1 平均 & 0.628 & 0.681 \\
逐图 F1 中位数 & 0.664 & 0.778 \\
逐图 F1 最小值 & 0.273 & 0.031 \\
逐图 F1 最大值 & 0.907 & 0.930 \\
Boundary IoU 平均 & 0.153 & 0.233 \\
\bottomrule
\end{tabular}

\vspace{6pt}
\begin{tabular}{lrrr}
\toprule
亮度分组 & 图像数 & 基线 F1 & 初版 F1 \\
\midrule
较低 & 11 & 56.56 & 52.82 \\
中间 & 11 & 61.90 & 74.88 \\
较高 & 10 & 69.43 & 83.00 \\
\bottomrule
\end{tabular}
\end{table}

32 张图中有 27 张 F1 提高、5 张下降。中位数从 0.664 升到 0.778，但最小值从 0.273 降到 0.031：多数图像受益，少数图像变差，而且变差得更厉害。逐图 F1 几乎呈两段分布：11 张不低于 0.85，5 张低于 0.2，0.2 到 0.5 之间没有样本。这个形态与“方法依赖一个占主导的肤色连通区域”的解释一致；样本只有 32 张，形态也可能带有偶然性。

按图像配对重采样 10,000 次，总体 F1 差值的 95\% 区间为 $+3.57$ 至 $+12.00$ 个百分点；逐图平均 F1 差值的区间为 $-0.68$ 至 $+10.60$ 个百分点。后者跨过 0，说明改善集中在部分图像上；只报一个总体 F1 会掩盖这一点。

按真值脸部平均 $Y$ 排序，把测试图分成 11、11、10 张的较低、中间、较高亮度三组。较低亮度组的 F1 反而从 56.56 降到 52.82，正确像素减少约 9.6 万、漏检增加约 9.6 万；中、高亮度组则分别提高 13 和 13.6 个百分点。低亮度下肤色色度偏离训练分布，$k=1.25$ 把真肤色也切掉了，减少的误检不足以补偿。需要说明，亮度分组是事后按真值区域平均 $Y$ 划分的，只有 11 张图，且亮度与场景、肤色、妆容混杂；这里只能说明低亮度条件下风险更高，不能把下降单独归因于亮度。

另外采用 Boundary IoU 观察轮廓贴合程度\cite{cheng2021boundary}：使用预测和真值各自的内侧边界带，带宽取图像对角线的 2\%（当前为 14 像素），逐图平均后从 0.153 升到 0.233。区域重合改善的同时，边界也有所改善，但 0.233 仍然不高；掩膜看起来完整，边缘并不可信。

\begin{figure}[H]
\centering
\includegraphics[width=\textwidth]{sample_stability.png}
\caption{逐图 F1 比较与累积分布。左图每个点对应一张测试图，虚线上方表示初版更高；右图为 F1 累积分布。}
\label{fig:stability}
\end{figure}

图~\ref{fig:stability} 同时给出逐图变化和低分尾部。5 张退步图中，有 3 张的 Precision 是 100\%——00609、00583、00735 的召回分别只有 1.6\%、4.2\%、7.4\%。掩膜里留下的像素确实都在脸上，但整张脸几乎没被找到。只看“干不干净”，它们接近满分；按 F1 算，它们和全黑没有本质区别。

\subsection{三类失败}

把 32 张误差图逐张看过去，失败可以归成三类。

第一类是颜色候选不足。00609 在 $k=1.25$ 下只剩 5,685 个候选像素，分成 56 个连通块，最大块 1,063 像素；最终召回 1.6\%，F1 0.031。图像整体偏冷，肤色被推到训练范围之外；错误在颜色阶段就发生了，后面的区域选择和填孔只是没能救回来。

\begin{figure}[H]
\centering
\includegraphics[width=\textwidth]{test_00609.png}
\caption{00609 的原图、真值、初版预测与误差对照。颜色候选不足导致大面积漏检。}
\label{fig:00609}
\end{figure}

第二类是区域选择错误。00718 中，棕色头发形成 58,716 像素的最大连通块，其中正确像素只有 2,705 个；填孔又增加 110 个正确像素和 10,871 个误检像素，最终 F1 0.045。面积只能表示区域大小，不能表示区域属于什么。

\begin{figure}[H]
\centering
\includegraphics[width=\textwidth]{test_00718.png}
\caption{00718 的原图、真值、初版预测与误差对照。头发被选为主要候选，填孔扩大了误检。}
\label{fig:00718}
\end{figure}

第三类不同：主体大致正确，但眼睛、嘴部与外部背景连通，填孔补不回来。00570 是一个清楚的例子，嘴部有一条横向缺口；它的 F1 是 0.759，召回 0.631，主要缺失不是“选错”，而是“补不上”。前两类属于颜色与区域的边界，在现有框架内没有便宜的修法；第三类只涉及“如何把开口封住”，可以单独做一次实验。下一章从这里开始。
